
% ════════════════════════════════════════════
% 模型准备（5.X.2.模型准备.tex）写作规则 —— 模板内嵌，AI 先读本文件再写
% ════════════════════════════════════════════
% 【定位】本块为每个问题的「模型准备」（三级子块），放在具体分析之后，用于把原始数据/问题转化为可建模样。
% 【结构】本块只有一个三级标题 \subsubsection{模型准备}；下面的「数据预处理」「算法选择」用加粗四级标题 \textbf{1.数据预处理} / \textbf{2.算法选择}（可带序号，无单独 \subsubsection）。
% 【数据预处理（加粗小标题 \textbf{数据预处理}，文字+公式+图）】① 开头引文；② 表达预处理原因；③ 讲预处理步骤及其对应方法；④ 理论结合公式表达预处理方法；⑤ 画预处理前后对比图；⑥ 编码结果展示（选）。
% 【算法选择（加粗小标题 \textbf{算法选择}，文字+表格）】① 开头过渡文字，承接预处理引入算法选择；② 引入对问题算法适用性讨论；③ 用表格展示算法对比结果；④ 对表格结果进行分析、表明要选的算法；⑤ 结尾段引出模型求解。
% 【表格】同 4.符号说明 的 longtable 规范（列宽总和=1.04-0.04*N；表头列名不加粗）。
% 【图片】预处理前后对比图可用同一 figure 内两张 minipage 并排；图不画 set_title()，标题由 caption 承担。
% ════════════════════════════════════════════
\subsubsection{模型准备}

\textbf{1.数据预处理}

% ① 开头引文 → ② 预处理原因 → ③ 步骤及方法 → ④ 公式表达 → ⑤ 前后对比图（⑥ 可选编码结果）
...（说明预处理原因）。为便于后续建模，首先对数据进行预处理，具体步骤如下。

将...采用...方法处理：设...为...，则预处理公式为：
\begin{equation}
    ... = ...
\end{equation}
其中，$...$表示...。

% 【图】在此插入「数据预处理前后对比图」（同一 figure 内两张 minipage 并排，宽高比档位定 width）。
% 模板不预置绘图方块，避免空白编译时因缺图而失败。图引用 \ref{fig:preprocessX}，label 为 fig:preprocessX。

\textbf{2.算法选择}

% ① 过渡承接 → ② 算法适用性讨论 → ③ 对比表 → ④ 分析→ ⑤ 引出模型求解
...（承接预处理）。针对本问，需在...等候选算法中择优，从...等维度对比其适用性，如表\ref{tab:algoX}所示。

% 算法对比表：列数 ≤5；方法名列常过长 → 行列转置（列=评估维度，行=方法名）
\begin{longtable}{>{\centering\arraybackslash}p{0.20\textwidth}>{\centering\arraybackslash}p{0.20\textwidth}>{\centering\arraybackslash}p{0.20\textwidth}>{\centering\arraybackslash}p{0.20\textwidth}}
  \caption{算法对比}
  \label{tab:algoX} \\
  \toprule
  算法 & 维度1 & 维度2 & 维度3 \\
  \midrule
  \endfirsthead
  \caption{算法对比（续）} \\
  \toprule
  算法 & 维度1 & 维度2 & 维度3 \\
  \midrule
  \endhead
  \bottomrule
  \endfoot
  ... & ... & ... & ... \\
  ... & ... & ... & ... \\
\end{longtable}

如\ref{tab:algoX}所示，...。综合以上分析，本问选择...算法进行求解。
